\subsection{Which training loss controls rejection?}
Fix a finite menu $\mathcal H$ of $M$ complete candidate services before validation. Each candidate specifies its proposal, rounding, certificate algorithm and fallback. On a context $x$, its first proposal is accepted when the valid complete-policy certificate $\Gamma_h(x)$ is at most $\tau>0$. For a one-date problem, \eqref{eq:gamma99} is such a certificate. At multiple dates it must additionally include the certified future-policy allowance. Define
\begin{equation}\label{eq:loss99}
 R_h(x)=\boldsymbol 1\{\Gamma_h(x)>\tau\},\qquad
 \ell_h(x)=\min\{1,\Gamma_h(x)/\tau\}.
\end{equation}
Then $R_h\le\ell_h$. Equality at the admission boundary is accepted. A valid fallback completes every rejected context; discarding it would change the service and invalidate its resource comparison.

\begin{theorem}[Prospective rejection and work guarantees]\label{thm:selection99}
Condition on all acquisition information that fixes $\mathcal H$. Let $X_1,\ldots,X_n$ be independent draws from a declared deployment-context law, independent of that information. For $0<\delta<1$, set
\[
 e_n=\sqrt{\frac{\log(2M/\delta)}{2n}}.
\]
An empirical minimizer $\widehat h$ of $R_h$ satisfies, with conditional probability at least $1-\delta$,
\begin{equation}\label{eq:rejection99}
 \mathbb E R_{\widehat h}(X)\le\min_{h\in\mathcal H}\mathbb E R_h(X)+2e_n.
\end{equation}
For an empirical minimizer of $\ell_h$, the same event instead gives
$\mathbb E R_{\widehat h}\le\min_h\mathbb E\ell_h+2e_n$.
If a verified resource bound is $W_h(x)\le C_0+C_1R_h(x)$, with common nonnegative constants, the rejection selector satisfies $\mathbb E W_{\widehat h}\le C_0+C_1(\min_h\mathbb E R_h+2e_n)$, and the surrogate selector has the analogous bound with $\ell_h$. More generally, directly minimizing empirical $W_h/W_{\max}\in[0,1]$ yields
\[
 \mathbb E W_{\widehat h}\le\min_h\mathbb E W_h+2W_{\max}e_n.
\]
All selection and rejected-candidate validation costs are additional one-time acquisition charges.
\end{theorem}

The work statement means $\mathbb E W_{\widehat h}\le C_0+C_1(\min_h\mathbb E R_h+2e_n)$ for the rejection selector, and analogously for $\ell$. A measured maximum is not a verified $W_{\max}$, and a rejection-optimal selector need not minimize runtime when fallback costs differ across contexts. On the finite dyadic domain used by the investment service, a finite fallback can supply a worst-case operation cap; its size, rather than a favorable observed clock, enters this theorem. Conditional concentration allows data-dependent acquisition, but does not allow the menu or certificate rule to be changed after looking at the validation sample.

Theorems~\ref{thm:action99}--\ref{thm:selection99} provide the missing distinction between optimizing a residual certificate and optimizing an action-relevant quantity. They do not retroactively certify the empirical performance of the R97 residual objective, or supply observations for the separately frozen R98 actors. Decision-aware prediction and surrogate calibration are established ideas \citep{ElmachtoubGrigas2022,LiuGrigas2021}. Here the loss is attached to an outward original-law, all-action Bellman certificate and a complete feasible service. The finite-menu concentration argument is classical; the support-to-admission construction specifies what that argument is permitted to select.
